\DocumentMetadata{
  pdfversion=2.0,
  pdfstandard=ua-2,
  lang=en-US,
  tagging=on,
  tagging-setup={math/setup=mathml-SE,tabsorder=structure}
}
\documentclass[11pt,letterpaper]{article}
\usepackage[margin=0.68in,includefoot]{geometry}
\usepackage{newcomputermodern}
\usepackage{amsmath,amscd,mathtools}
\usepackage{tikz,adjustbox}
\providecommand{\square}{\mdlgwhtsquare}
\usepackage[hidelinks]{hyperref}
\hypersetup{
  pdftitle={Numerical Analysis PhD Exam, May 1995},
  pdfauthor={University of Florida Department of Mathematics},
  pdfsubject={Graduate mathematics examination},
  pdfdisplaydoctitle=true,
  bookmarksopen=true
}
\setcounter{secnumdepth}{0}
\setlength{\parindent}{0pt}
\setlength{\parskip}{0.28em}
\setlength{\emergencystretch}{3em}
\setlength{\leftmargini}{2.15em}
\setlength{\leftmarginii}{2.65em}
\setlength{\leftmarginiii}{3em}
\pagestyle{empty}
\raggedbottom
\allowdisplaybreaks
\NewTaggingSocketPlug{block/list/label}{examproblem}{%
  \tagstructbegin{tag=Lbl}\tagstructend%
  \tagstructbegin{tag=\LBody}%
  \tagstructbegin{tag=\ProblemHeadingTag,title-o={\ProblemHeadingText}}%
  \tagmcbegin{tag=\ProblemHeadingTag,actualtext={\ProblemHeadingText}}%
  #2%
  \tagmcend%
  \tagstructend%
}
\newenvironment{examproblems}[2]{%
  \def\ProblemHeadingTag{#2}%
  \AssignTaggingSocketPlug{block/list/label}{examproblem}%
  \begin{enumerate}%
  \setcounter{enumi}{#1}%
  \renewcommand{\labelenumi}{\textbf{\theenumi.}}%
  \setlength{\topsep}{0.35em}%
  \setlength{\partopsep}{0pt}%
  \setlength{\itemsep}{0.48em}%
  \setlength{\parsep}{0.18em}%
}{\end{enumerate}}
\newenvironment{examsubparts}{%
  \AssignTaggingSocketPlug{block/list/label}{default}%
  \begin{enumerate}%
  \setlength{\topsep}{0.18em}%
  \setlength{\partopsep}{0pt}%
  \setlength{\itemsep}{0.16em}%
  \setlength{\parsep}{0.1em}%
}{\end{enumerate}}
\newenvironment{examinstructions}{%
  \par\begingroup\itshape
}{%
  \par\endgroup\vspace{0.18em}
}
\makeatletter
\renewcommand\subsection{\@startsection{subsection}{2}{0pt}%
  {-1.25ex plus -.25ex minus -.1ex}%
  {0.55ex plus .1ex}%
  {\normalfont\large\bfseries}}
\renewcommand{\@maketitle}{%
  \newpage\null\vskip -1em%
  {\footnotesize\bfseries
    \href{https://www.ufl.edu/}{University of Florida}, \href{https://math.ufl.edu/}{Department of Mathematics}\par}%
  \vskip -0.5em%
  \tagstructbegin{tag=H1,title={Numerical Analysis PhD Exam, May 1995}}%
  \tagpdfparaOff%
  \tagmcbegin{tag=H1}%
    {\LARGE\bfseries\href{https://gma.math.ufl.edu/exams/numerical-analysis-phd/}{Numerical Analysis PhD Exam}, May 1995\par}%
  \tagmcend%
  \tagpdfparaOn%
  \tagstructend%
  \vskip 0.25em%
  \tagmcbegin{artifact}%
    \hrule height 1pt%
  \tagmcend%
  \vskip 1em%
}
\makeatother
\title{Numerical Analysis PhD Exam, May 1995}
\author{}
\date{}
\begin{document}
\maketitle
\thispagestyle{empty}
\begin{examproblems}{0}{H2}
\def\ProblemHeadingText{Problem 1.}
\item
Consider the linear system \(Ax=b\).
\begin{examsubparts}
\item[\textnormal{(a)}]
Suppose the columns of~\(A\) are linearly independent. Derive the normal equation for the~\(x\) that minimizes the 2-norm of the residual \(r=b-Ax\).
\item[\textnormal{(b)}]
Suppose that the rows of~\(A\) are linearly independent. Derive the normal equation for the solution to \(Ax=b\) that has the smallest 2-norm.
\item[\textnormal{(c)}]
Using the singular value decomposition, obtain a formula for the least squares solution to \(Ax=b\). That is, if~\(M\) denotes the set of all~\(x\) for which the 2-norm of the residual is minimal, your formula should give the point in~\(M\) with smallest 2-norm.
\end{examsubparts}
\def\ProblemHeadingText{Problem 2.}
\item
Suppose the columns of~\(A\) are independent. Given the QR factorization of~\(A\), give a formula for the distance from some given vector~\(b\) to the space spanned by the columns of~\(A\).
\def\ProblemHeadingText{Problem 3.}
\item
Verify the Woodbury formula
\[
(B-UV^T)^{-1}=B^{-1}+B^{-1}U(I-V^TB^{-1}U)^{-1}V^TB^{-1}.
\]
\def\ProblemHeadingText{Problem 4.}
\item
Explain how the singular value decomposition can be used to determine the rank of a matrix in the presence of rounding errors.
\def\ProblemHeadingText{Problem 5.}
\item
Derive a formula for the 2-norm of a matrix in terms of its singular values.
\def\ProblemHeadingText{Problem 6.}
\item
Let~\(f\) be an \((n+1)\)-times continuously differentiable function on the interval \([a,b]\), and let \(a\leq x_0<x_1<\cdots<x_n\leq b\). Let \(P_n\) be the unique polynomial interpolant of degree at most~\(n\), interpolating the data \((x_k,f(x_k))\), \(k=0,1,\ldots,n\). Prove that there exists a point \(\xi\in[a,b]\) such that
\[
f(x)-P_n(x)=\frac{f^{(n+1)}(\xi)}{(n+1)!}\prod_{k=0}^n(x-x_k)
\]
 for all \(x\in[a,b]\).
\def\ProblemHeadingText{Problem 7.}
\item
Let~\(f\) be a twice continuously differentiable function on the interval \([a,b]\), and \(x_0,x_1,\ldots,x_N\) be an equi-spaced partition of \([a,b]\), with step size~\(h\) and \(x_0=a\). Assume that only the values \(f(x_k)\), \(k=0,1,\ldots,N\), are known.
\begin{examsubparts}
\item[\textnormal{(a)}]
State the trapezoidal rule \(T(f,N)\) to approximate \(\int_a^b f(x)\,dx\).
\item[\textnormal{(b)}]
Prove that
\[
\int_a^b f(x)\,dx-T(f,N)=-\frac{(b-a)^3}{12N^2}f''(\xi)
\]
 for some point \(\xi\in(a,b)\).
\end{examsubparts}
\def\ProblemHeadingText{Problem 8.}
\item
Let~\(f\) be a twice continuously differentiable, \(2\pi\)-periodic, function on~\(\mathbb{R}\) and consider the problem of computing the Fourier coefficients
\[
\tag{8.1} c_n=\frac{1}{2\pi}\int_0^{2\pi}e^{-inx}f(x)\,dx
\]
 by using the values \(f_k=f(2k\pi/N)\) for \(k=0,1,\ldots,N-1\).
\begin{examsubparts}
\item[\textnormal{(a)}]
Define the discrete Fourier transform \(\{\hat f_n\}\) of the sequence \(\{f_k\}\).
\item[\textnormal{(b)}]
By applying the trapezoidal rule to the integral in (8.1), show that \(c_n\) can be approximated by \(\hat f_n\).
\item[\textnormal{(c)}]
By performing an analysis of the error in the approximation in (b), show that \(\hat f_n\) is a poor approximation of \(c_n\). Hint: Investigate behavior of error as \(N\to\infty\).
\item[\textnormal{(d)}]
Outline a method of obtaining a better approximation \(\tilde c_n\) of \(c_n\). (You do not need to actually find a formula for this \(\tilde c_n\).) For the \(\tilde c_n\) obtained by your method, define
\[
\lVert\tilde c-c\rVert_\infty=\max\{|\tilde c_n-c_n|,\ n=1,2,\ldots\}.
\]
 Find the order of convergence of \(\lVert\tilde c-c\rVert_\infty\) as \(N\to\infty\).
\end{examsubparts}
\def\ProblemHeadingText{Problem 9.}
\item
Give a careful statement of the contraction mapping theorem for functions defined on subsets of \(\mathbb{R}^n\) into \(\mathbb{R}^n\).
\begin{examsubparts}
\item[\textnormal{(a)}]
Since it is usually not easy to demonstrate the contraction property in particular examples, state a stronger condition which implies the contraction property.
\item[\textnormal{(b)}]
Discuss the relative merits of determining the fixed points of a function \(\varphi:\mathbb{R}\to\mathbb{R}\) by the iterations:
\[
\text{(FP)}\qquad x_{k+1}=\varphi(x_k)
\]

\[
\text{(DFP)}\qquad x_{k+1}=\varphi(\varphi(x_k))
\]

\[
\text{(STEF)}\qquad x_{k+1}=x_k-\frac{(\varphi(x_k)-x_k)^2}{\varphi(\varphi(x_k))-2\varphi(x_k)+x_k}.
\]
\end{examsubparts}
\def\ProblemHeadingText{Problem 10.}
\item
Let~\(A\) be a real, positive definite \(n\times n\) matrix and~\(b\) be an \(n\times1\) vector.
\begin{examsubparts}
\item[\textnormal{(a)}]
Derive the steepest descent algorithm for approximating the true solution~\(x\) of \(Ax=b\), by a sequence of iterates \(\{x_k\}\).
\item[\textnormal{(b)}]
Define the real-valued function~\(h\) on \(\mathbb{R}^n\) by
\[
h(x)=(Ax-b)^TA^{-1}(Ax-b)
\]
 and show that
\[
h(x_{k+1})\leq(1-c(A)^{-1})^2h(x_k)
\]
 where \(c(A)\) is the condition number of the matrix~\(A\).
\end{examsubparts}
\end{examproblems}
\end{document}
